%! AP Statistics — Unit 1, Lesson 1.4 — Warm-Up
\documentclass[10pt]{article}
\usepackage{apstats-article}
\usepackage{apstats-boxes}

%=============================================================================
\begin{document}

\pageheader{Unit 1, Lesson 1.4}{Warm-Up}
\namedateperiod

\vspace{0.18in}

\noindent In Lesson 1.3 you built frequency and relative frequency tables.
Today you will display the same data \textit{visually}. Use the table below
(from Lesson 1.3 guided practice) to answer the questions.

\vspace{0.1in}

\begin{center}
\small
\renewcommand{\arraystretch}{1.8}
\begin{tabularx}{0.72\linewidth}{@{} >{\bfseries}p{0.22\linewidth}
                                        C{0.2\linewidth}
                                        C{0.26\linewidth} @{}}
\rowcolor{sky}
\textbf{Mode} & \textbf{Frequency} & \textbf{Rel.\ Frequency} \\
\midrule
Car  & 10 & 0.40 \\
Bus  & 8  & 0.32 \\
Walk & 4  & 0.16 \\
Bike & 3  & 0.12 \\
\midrule
\rowcolor{sky}\textbf{Total} & 25 & 1.00 \\
\end{tabularx}
\end{center}

\vspace{0.18in}

\begin{enumerate}

  \item What is the most common mode of transportation?
        What percentage of students use it?\\[6pt]
        \writeline

  \item Without drawing anything yet --- if you wanted to show this data
        \textit{visually} so a reader could see at a glance which mode is most
        popular, what type of picture or shape would you draw?\\[6pt]
        \writeline\\[1pt]\writeline

  \item In the space below, make a quick \textbf{rough sketch} of how you
        would display this data. Label whatever you can think of.

\vspace{0.12in}
\begin{tcolorbox}[colback=white, colframe=navy, boxrule=0.8pt, arc=2mm,
    left=3mm, right=3mm, top=2mm, bottom=2mm, height=2.4in]
\end{tcolorbox}

  \item What information does your sketch show that the table does \textit{not}
        show as easily? What does the table show that your sketch might not?\\[6pt]
        \writeline\\[1pt]\writeline

\end{enumerate}

\vspace{0.12in}

\begin{tcolorbox}[colback=greenbg, colframe=greenacc, boxrule=0.8pt, arc=2mm,
    left=4mm, right=4mm, top=3mm, bottom=3mm]
\small
\begin{itemize}[leftmargin=1.3em, itemsep=5pt]
  \item Today we formalize two visual displays for categorical data:
        \textbf{bar charts} and \textbf{pie charts}.
  \item Both displays use the same relative frequencies from Lesson 1.3 ---
        but they communicate the distribution differently.
  \item We will also learn to spot graphs that \textit{lie} --- even when the
        underlying data are correct.
\end{itemize}
\end{tcolorbox}

\vspace{0.12in}

\noindent\textcolor{slate}{\small One question you have going into today's lesson,
or something you're curious about:}\\[8pt]
\writeline

\end{document}
